% pptfig/equations.tex —— 可直接粘进 MathType / Word 公式编辑器的 LaTeX 片段（PPT 用）
% 用法：MathType 里 Edit > Preferences > Workspace Preferences 勾选
%       "Allow TeX language input"，然后粘贴下面任意一段，回车即转为公式。
% 每段都是独立完整的 LaTeX 表达式（不含 \begin{equation}），方便逐条粘贴。

% (1) 灰箱 ESO：类 fal 校正（默认档 a=0.75, δ=0.10, λ=10）
\dot z_1 = z_2 + b_0 u + v + \beta_1 G(e_1;1,\delta,\lambda),\quad
\dot z_2 = \beta_2 G(e_1;a,\delta,\lambda),\quad
\beta_1 = 2\omega_o,\ \beta_2 = \omega_o^2

% (2) 类 fal 增益定义（z1 通道指数 1，z2 通道指数 a）
G(e;a,\delta,\lambda)=\left(|e|+\delta\lambda^{-1/a}\right)^{a}\tanh\left(\frac{\lambda e}{\delta}\right)

% (3) 控制律
u=\frac{u_0-z_2-v}{b_0},\quad u_0=k_p(r-y),\quad e_1=y-z_1

% (4) 低通重复通道（内模）
\tau\dot v+v=Q\,v(t-T)+\beta_3\tanh\left(\frac{e_1}{\zeta}\right)=:\chi,\quad e_f:=\chi-v=\tau\dot v

% (5) 记忆失配检测量（方向限定，用于反相/撤除事件）
D_m^{\downarrow}=\frac{\max\left(0,\,E_v(t-T)-E_v(t)\right)}{E_v(t-T)+\varepsilon_v},\quad
E_v(t)=\frac{1}{T}\int_{t-T}^{t}v^2(\tau)\,\mathrm{d}\tau

% (6) 圆盘质量不对中模式：角度锁相周期负载（机械角 φ，dφ/dt=ω）
% 1–6 次谐波，按给定顺序（3 次 0.10、4 次 0.45，非单调）；THD=77.46%（超过旧 30% 约束）
T(\varphi)=T(\varphi+2\pi),\quad
T_L=T_0+A\,\tilde T(\varphi),\quad
\tilde T(\varphi)=\frac{\sum_{n=1}^{6}a_n\sin n\varphi}{\max\left|\sum_{n=1}^{6}a_n\sin n\varphi\right|},\quad
(a_1,\ldots,a_6)=(1,\,0.55,\,0.10,\,0.45,\,0.25,\,0.15),\quad
\mathrm{THD}=\frac{\sqrt{\sum_{n=2}^{6}a_n^2}}{a_1}=77.46\%,\quad
\frac{\mathrm{d}\varphi}{\mathrm{d}t}=\omega

% (7) 小信号斜率与等效观测带宽
G'(0)=\delta^{a-1},\qquad \omega_o^{\mathrm{eff}}=\omega_o\,\delta^{(a-1)/2}
